% concatenated from Frequency.tex
\documentclass[11pt]{amsart}

\usepackage{amsmath,amssymb,latexsym} % ams auxiliaries.
\usepackage{graphicx}
\usepackage{fullpage}
\usepackage[square, numbers]{natbib}   % Arguably a nicer bibliography.
\usepackage{comment}


% \usepackage{times}   % Causes problems with accents & spacing in PDF.





\newtheorem{theorem}{Theorem} % If subsection, get 3.1.2, etc.
\newtheorem{corollary}{Corollary}
\newtheorem{proposition}{Proposition}
\newtheorem{claim}{Claim}
\newtheorem{lemma}{Lemma}
\theoremstyle{definition}
\newtheorem{definition}{Definition}
\newtheorem{example}{Example}
\newtheorem{question}{Question}
\newtheorem{conjecture}{Conjecture}
\newtheorem{remark}{Remark}



\def\Ker{\operatorname{Ker}}
\def\M{\mathsf{M}}
\def\K{\mathsf{K}}
\def\R{\mathbb{R}}
\def\Z{\mathbb{Z}}
\def\T{\mathsf{T}}
\def\F{\mathcal{F}}
\def\aa{\boldsymbol{a}}
\def\bb{\boldsymbol{b}}
\def\cc{\boldsymbol{c}}
\def\ee{\boldsymbol{e}}
\def\qq{\boldsymbol{q}}
\def\rr{\boldsymbol{r}}
\def\ss{\boldsymbol{s}}
\def\uu{\boldsymbol{u}}
\def\vv{\boldsymbol{v}}
\def\xx{\boldsymbol{x}}
\def\yy{\boldsymbol{y}}
\def\zz{\boldsymbol{z}}
\def\llambda{\boldsymbol{\lambda}}
\def\aalpha{\boldsymbol{\alpha}}

\def\zero{\boldsymbol{0}}
\def\supp{\operatorname{supp}}
\def\conv{\operatorname{conv}}
\def\ds{\displaystyle}
\def\KG{\operatorname{KG}}
\def\cd{\operatorname{cd}}
\def\rk{\operatorname{rk}}
\def\alt{\operatorname{alt}}
%\def\H{\mathcal{H}}
\def\id{\operatorname{id}}
\def\diam{\operatorname{diam}}
\def\stab{\operatorname{stab}}
\def\sd{\operatorname{sd}}
\def\mod{\operatorname{mod}}
\def\coind{\operatorname{coind}}

\title{Lower bound for the tau-matching chromatic number via the equitable colorability defect}
\author{Saeed Shaebani, Meysam Alishahi}
\date{arxiv:2101.04701v2; CC BY 4.0}
\begin{document}
\maketitle
\noindent\textit{Editorial note.} The counted unit is the article's own Theorem mainthmSaeed2 at source lines 301-309 together with the proof environment that begins on the next line (311-345); the pairing is the ordinary adjacent one, with no other proof environment in between.\par\medskip
\noindent\textit{Editorial note.} The retained context at source lines 115-121, 124-126, 264-283 and 294-299 is the article's own definition of the tau-matching chromatic number, of the general Kneser hypergraph and of the equitable colorability defect, together with the cited lower bound for the Kneser chromatic number that the proof applies and the announcement of the theorem.\par\medskip
\noindent\textit{Editorial note.} The native preamble is retained from source lines 1-3 and 5-69; only the hyperref line at source line 4 is omitted, so that the material carries no page-backref and no page-label machinery, and no other preamble line is altered.\par\medskip
\noindent\textit{Editorial note.} The bibliography fragment contains verbatim only the five bibliography entries that the retained lines cite; no citation in the retained text is left without an entry.\par\medskip
\noindent\textit{Editorial note.} The same source's other theorem, mainthmSaeed, and its sharpness Proposition are not counted here and are not used as premises.\par\medskip
\noindent\textit{Editorial scope.} One original result: the article's Theorem mainthmSaeed2 with its complete author proof, plus the article's own definitions of the tau-matching chromatic number, of the general Kneser hypergraph and of the equitable colorability defect. The article's Theorem mainthmSaeed, its sharpness Proposition and the already counted public unit from this source are not counted here and are not used as premises.\par\medskip
\noindent\textit{Allowed standard prerequisites (editorial).}\begin{itemize}
\item General Kneser hypergraph and its chromatic number. For a hypergraph H the general Kneser hypergraph is the hypergraph whose vertices are the edges of H and whose edges are the r-tuples of pairwise disjoint edges of H; the article's own definition is retained at source lines 124-126, together with the remark there that its chromatic number is the special case of the tau-matching chromatic number in which every frequency equals r-1.. The proof colours this hypergraph with the colour set A and applies the lower bound for its chromatic number at source lines 324-326.
\item Equitable colorability defect and the cited lower bound. The r-th equitable colorability defect of H is the minimum number of vertices whose removal leaves a subhypergraph admitting a proper vertex colouring with r colours whose class sizes differ by at most one; the article states at source lines 275-279 that the chromatic number of the general Kneser hypergraph of H is at least the ceiling of that defect divided by r-1. This is a cited result of Abyazi Sani, Alishahi, Azarpendar and Jafari, not a claim of the counted theorem.. Applied to the augmented hypergraph at source lines 324-326 to convert the number of colours into a bound on the defect.
\item Admissible colour frequencies and matching colourings. An r-colour-frequency mapping assigns to each colour an integer below r; an (A,tau)-matching colouring of H colours the edges by colours of A so that every nonempty monochromatic matching has size at most the frequency of its colour; H is tau-admissible when such a finite A exists, and the tau-matching chromatic number is the least size of such an A (source lines 115-121).. Fixes the two objects and both conclusions of the counted theorem.
\end{itemize}\medskip
Let ${\mathcal H}$ be a hypergraph with $n$ vertices, and $r$ be a positive integer. Also, let $C$ be an arbitrary nonempty (finite or infinite) set, and regard its elements as our available colors.

\noindent Any arbitrary mapping $\tau: C\longrightarrow \{0,1,\ldots,r-1\}$ is called an {\it $r$-color-frequency mapping}. So, by an $r$-color-frequency mapping, each color in $C$ receives a nonnegative integer less that $r$ as its {\it frequency}. Our aim is coloring the edges of ${\mathcal H}$ by some colors in $C$ in such a way that the size of each nonempty monochromatic matching does not exceed the frequency of the color associated to its edges. Naturally, one is looking for the least possible number of colors needed for such a coloring.

\noindent Formally, for a subset $A$ of $C$, an {\it $(A,\tau)$-matching coloring} of ${\mathcal H}$ is
a coloring of edges of ${\mathcal H}$ by some colors in $A$ in such a way that the size of each nonempty monochromatic matching is bounded above by the frequency of the corresponding color associated to its edges.
The hypergraph ${\mathcal H}$ is called {\it $\tau$-admissible} if ${\mathcal H}$ admits an $(A,\tau)$-matching coloring for some finite subset $A$ of $C$; and in this case, we define the {\it $\tau$-matching chromatic number} of ${\mathcal H}$, denoted by $\chi_M(\tau, {\mathcal H})$, equal to the minimum cardinality of such a finite subset $A$ of $C$. As usual, for the case where ${\mathcal H}$ is not $\tau$-admissible, we set $\chi_M(\tau, {\mathcal H}):=+\infty$.

The {\it general Kneser hypergraph} ${\rm KG}^r({\mathcal H})$ is an $r$-uniform hypergraph which has $E({\mathcal H})$ as its vertex set and  whose edges are formed by $r$ pairwise disjoint edges of ${\mathcal H}$, that is,
$$E\left({\rm KG}^r({\mathcal H})\right) :=\left\{\{e_1,\ldots,e_r\}:\; e_i\cap e_j=\varnothing \mbox{ for all } i\neq j\right\}.$$
It is worth pointing out that the well-known and extensively well-studied chromatic number of general Kneser hypergraph $\chi \left( {\rm KG}^r({\mathcal H}) \right)$ is a special case of $\chi_M(\tau, {\mathcal H})$ where all color frequencies are the fixed constant $r-1$.

\def\cprime{$'$} \def\cprime{$'$}
For two positive integers $n$ and $k$, the symbol ${[n]\choose k}$  denotes the set of all $k$-subsets of $[n]$.
The {\it complete $k$-uniform  hypergraph}  $K_n^k$ is a hypergraph with vertex set $[n]$ and the edge set ${[n]\choose k}$.

As a generalization of these results, Dol{\cprime}nikov and K{\v{r}}{\'{\i}}{\v{z}} \cite{MR953021,MR1081939,MR1665335}
proved that for all hypergraphs ${\mathcal H}$, the relation
$$\chi\left({\rm KG}^r   ({\mathcal H})  \right) \geq \left\lceil \frac{ {\rm cd}^r({\mathcal H}) }{r-1} \right\rceil $$
holds, where the symbol ${\rm cd}^r({\mathcal H})$, called the {\it $r$-th colorability defect} of ${\mathcal H}$, is the least possible number of vertices that one must remove from
${\mathcal H}$ in such a way that the vertices of the remaining subhypergraph could be properly colored by at most $r$ colors. This lower bound for
$\chi\left({\rm KG}^r   ({\mathcal H})  \right) $ is sharp since for
$n \geq r(k-1) + 1$ we have
${\rm cd}^r \left(  K_n^k \right) = n-r(k-1)$; and therefore,
$$\begin{array}{lrr}
	\chi\left({\rm KG}^r \left(  K_n^k \right) \right) = \left\lceil \frac{{\rm cd}^r \left(  K_n^k \right)}{r-1} \right\rceil &  {\rm if }  &  n \geq r(k-1) + 1 .
\end{array}$$
Later, Abyazi Sani and Alishahi \cite{ABYAZISANI2018229}, and Azarpendar and Jafari \cite{AZARPENDAR2023103664},
refined Dol{\cprime}nikov and K{\v{r}}{\'{\i}}{\v{z}}'s Theorem by showing that
$$\chi\left({\rm KG}^r ({\mathcal H}) \right) \geq \left\lceil \frac{ {\rm ecd}^r({\mathcal H}) }{r-1} \right\rceil ,$$
where ${\rm ecd}^r({\mathcal H})$, referred to as the {\it $r$-th equitable colorability defect} of ${\mathcal H}$, denotes the minimum number of vertices that must be
removed from ${\mathcal H}$ in order to obtain a remaining subhypergraph that admits a proper vertex coloring using $r$ colors, such that the sizes of the color classes differ by at most one.
This inequality is also sharp because
$$\begin{array}{lrr}
	\chi\left({\rm KG}^r \left(  K_n^k \right) \right) = \left\lceil \frac{{\rm ecd}^r \left(  K_n^k \right)}{r-1} \right\rceil &  {\rm whenever }  &  n \geq r(k-1) + 1 .
\end{array}$$

In this section, we are concerned with providing a tight lower bound for $\chi_M(\tau, {\mathcal H})$ in terms of equitable colorability defect.
In this regard, we state and prove the following general Theorem \ref{mainthmSaeed2}.
As a specific instance of the assertion of Theorem \ref{mainthmSaeed2}, the celebrated inequality
$\chi\left({\rm KG}^r ({\mathcal H}) \right) \geq \left\lceil \frac{ {\rm ecd}^r({\mathcal H}) }{r-1} \right\rceil $
follows by considering the fixed constant $r$-color-frequency mapping
$\tau: \mathbb{N} \longrightarrow \{0,1,\ldots,r-1\}$ with $\tau (a) = r-1$ for all elements $a$ in $ \mathbb{N} $.

\begin{theorem} \label{mainthmSaeed2}
	Let $C$ be a nonempty finite or infinite set, and $r$ be an integer with $r\geq 2$.
	Also, let the function $\tau: C \longrightarrow \{0,1,\ldots,r-1\}$ be an $r$-color-frequency mapping, and ${\mathcal H}$ be a $\tau$-admissible hypergraph. Then
	there exists a finite subset
	$B$ of $C$ for which
	$$\sum_{b\in B}\tau(b) \geq  {\rm ecd}^r({\mathcal H}) .$$
	Furthermore,
	$$\chi_M(\tau,{\mathcal H})\geq \min\left\{|A|:\   A \subseteq C, \mbox{ and } A \mbox{ is finite,}  \mbox{ and also } \sum_{a\in A}\tau(a) \geq  {\rm ecd}^r({\mathcal H})\right\}.$$
\end{theorem}

\begin{proof}{
		The main idea of the proof is to show that if ${\mathcal H}$ admits an $(A,\tau)$-matching coloring for some finite subset $A$ of $C$, then we would have $\displaystyle\sum_{a\in A}\tau(a) \geq  {\rm ecd}^r({\mathcal H})$.
		
		Let us regard an arbitrary finite subset $A$ of $C$ for which ${\mathcal H}$ admits an $(A,\tau)$-matching coloring $c: E({\mathcal H})\longrightarrow A$. We construct a new hypergraph $\widetilde{{\mathcal H}}$ in such a way that for each element $a$ of $A$ with $\tau(a) < r-1$, we add $r-1-\tau(a)$ new vertices
		$v^{(a)}_{1},v^{(a)}_{2},\ldots ,v^{(a)}_{r-1-\tau(a)}$ to the vertex set of ${\mathcal H}$ and also we add all new singleton edges $\left\{v^{(a)}_{1}\right\},\left\{v^{(a)}_{2}\right\},\ldots ,\left\{v^{(a)}_{r-1-\tau(a)}\right\}$ to the edge set of ${\mathcal H}$. Denote the resulting hypergraph by $\widetilde{{\mathcal H}}$. Now, extend the coloring $c: E({\mathcal H})\longrightarrow A$ to a mapping $\widetilde{c}: E \left(\widetilde{{\mathcal H}}\right)\longrightarrow A$ with the property that for each element $a$ in $A$ with $\tau(a) < r-1$ we have
		$$\widetilde{c}\left(\left\{v^{(a)}_{1}\right\}\right)=\widetilde{c}\left(\left\{v^{(a)}_{2}\right\}\right)=\cdots =\widetilde{c}\left(\left\{v^{(a)}_{r-1-\tau(a)}\right\}\right)=a.$$
		We claim that $\widetilde{c}$ is a proper vertex coloring of ${\rm KG}^r\left(\widetilde{{\mathcal H}}\right)$.
		Suppose the assertion of this claim is false. Then, there exist some $a$ in $A$ and an edge of ${\rm KG}^r\left(\widetilde{{\mathcal H}}\right)$, say
		$\{e_{1},e_{2},\ldots , e_{r}\}$, in such a way that $\widetilde{c}(e_{1})=\widetilde{c}(e_{2})=\cdots = \widetilde{c}(e_{r}) =a$ holds.
		Since $|\{e_{1},e_{2},\ldots , e_{r}\}\setminus E(\mathcal{H})|\leq r-1-\tau (a)$,
		we obtain the existence of $\tau (a) +1$ pairwise disjoint edges of $\mathcal{H}$ that are colored by $a$ under $c: E({\mathcal H})\longrightarrow A$, contradicting the fact that $c$ is an $(A,\tau)$-matching coloring of $\mathcal{H}$.
		
		From what has already been proved, it may be concluded that 
		$$|A| \geq \chi\left({\rm KG}^r\left(\widetilde{{\mathcal H}}\right)\right)\geq { {\rm ecd}^r\left(\widetilde{{\mathcal H}}\right)\over r-1};$$
		and consequently, it follows that 
		$$(r-1)|A| \geq {\rm ecd}^r\left(\widetilde{{\mathcal H}}\right).$$
		
		We recall that ${\rm ecd}^r\left(\widetilde{{\mathcal H}}\right)$ denotes the minimum number of vertices that must be
		removed from $\widetilde{{\mathcal H}}$ in order to obtain a remaining {\it equitably $r$-colorable} subhypergraph, that is, a hypergraph which admits a proper vertex coloring using $r$ colors, such that the sizes of the color classes differ by at most one.
		
		\noindent
		On one hand, if a hypergraph contains a singleton edge, it is never $r$-colorable. Therefore, in order to obtain an induced subhypergraph of $\widetilde{\mathcal{H}}$ that is equitably $r$-colorable, it is necessary to remove all vertices in $V(\widetilde{\mathcal{H}}) \setminus V(\mathcal{H})$. On the other hand, eliminating singleton edges alone
		may be insufficient; at least ${\rm ecd}^r(\mathcal{H})$ vertices from $\mathcal{H}$ must also be removed to acquire a remaining subhypergraph that is equitably $r$-colorable.
		Thus, 
		$${\rm ecd}^r\left(\widetilde{{\mathcal H}}\right) \geq \left|V\left(\widetilde{{\mathcal H}}\right)\right|-|V({\mathcal H})|+{\rm ecd}^r({\mathcal H}).$$
		
		Now, since 
		$(r-1)|A| = \left( \displaystyle\sum_{a\in A}\tau(a) \right) + \left|V\left(\widetilde{{\mathcal H}}\right)\right|-|V({\mathcal H})|$,
		the inequality
		$(r-1)|A| \geq {\rm ecd}^r\left(\widetilde{{\mathcal H}}\right)$
		implies
		$\displaystyle\sum_{a\in A}\tau(a) \geq  {\rm ecd}^r({\mathcal H})$;
		which is the desired conclusion.
	}
\end{proof}

\begin{thebibliography}{99}
\bibitem{ABYAZISANI2018229}
R.~{Abyazi Sani} and M.~Alishahi.
\newblock A new lower bound for the chromatic number of general {K}neser
  hypergraphs.
\newblock {\em European Journal of Combinatorics}, 71:229--245, 2018.

\bibitem{AZARPENDAR2023103664}
S.~Azarpendar and A.~Jafari.
\newblock Lower bounds for the chromatic number of certain {K}neser-type
  hypergraphs.
\newblock {\em European Journal of Combinatorics}, 110:103664, 2023.

\bibitem{MR953021}
V.~L. Dol{\cprime}nikov.
\newblock A combinatorial inequality.
\newblock {\em Sibirsk. Mat. Zh.}, 29(3):53--58, 219, 1988.

\bibitem{MR1081939}
I.~K{\v{r}}{\'{\i}}{\v{z}}.
\newblock Equivariant cohomology and lower bounds for chromatic numbers.
\newblock {\em Transactions of the American Mathematical Society},
  333(2):567--577, 1992.

\bibitem{MR1665335}
I.~K{\v{r}}{\'{\i}}{\v{z}}.
\newblock A correction to: ``{E}quivariant cohomology and lower bounds for
  chromatic numbers'' [{T}ransactions of the {A}merican {M}athematical
  {S}ociety, {333} (1992), no. 2, 567--577; {MR}1081939 (92m:05085)].
\newblock {\em Transactions of the American Mathematical Society},
  352(4):1951--1952, 2000.
\end{thebibliography}
\end{document}
